\section{Cơ chế tự phục hồi đa tầng}
\label{sec:self-recovery}

Mã nguồn do LLM sinh ra có thể chứa lỗi cú pháp, lỗi ngữ nghĩa, hoặc vi phạm ràng buộc dữ liệu. Reflexion~\cite{shinn2023reflexion} và MetaGPT~\cite{hong2023metagpt} đã cho thấy phản hồi có cấu trúc giúp LLM sửa lỗi hiệu quả hơn so với thử lại mà không cung cấp ngữ cảnh. \tool áp dụng cơ chế tự phục hồi đa tầng hoạt động ở ba cấp độ, mỗi cấp có phạm vi và chiến lược phục hồi khác nhau.

\subsection{Phân loại lỗi}

Khi một tác tử tạo ra ngoại lệ, hệ thống phân loại lỗi dựa trên kiểu ngoại lệ và từ khóa trong thông điệp. Bảng~\ref{tab:error-taxonomy} mô tả ba loại lỗi: \textit{transient} (lỗi tạm thời do mạng hoặc giới hạn API, có thể tự phục hồi khi thử lại), \textit{recoverable} (lỗi sửa được bằng phản hồi có cấu trúc, ví dụ lỗi cú pháp), và \textit{permanent} (lỗi logic nghiêm trọng không thuộc hai loại trên). Phân loại này phục vụ quyết định phục hồi có chọn lọc ở các tầng tiếp theo.

\begin{table}[htbp]
\centering
\caption{Phân loại lỗi tự động theo mức độ nghiêm trọng.}
\label{tab:error-taxonomy}
\begin{tabular}{lp{4cm}p{6cm}}
\toprule
\textbf{Loại} & \textbf{Điều kiện} & \textbf{Mô tả} \\
\midrule
Transient   & Timeout, lỗi kết nối, rate limit & Lỗi tạm thời --- có khả năng tự phục hồi khi thử lại \\
Recoverable & Lỗi cú pháp, giá trị không hợp lệ & Lỗi sửa được bằng phản hồi có cấu trúc cho LLM \\
Permanent   & (mặc định) & Lỗi logic hoặc lỗi hệ thống không tự khắc phục được \\
\bottomrule
\end{tabular}
\end{table}

\subsection{Tầng 1: Cổng xác thực và thử lại tại tác tử sinh mã}

Các tác tử sinh mã ($\mathcal{A}_{builder}$, $\mathcal{A}_{task}$) có vòng lặp tự kiểm tra: sau mỗi lần gọi LLM, tác tử kiểm tra sự tồn tại của sản phẩm mục tiêu. Nếu sản phẩm chưa được tạo, tác tử bổ sung hướng dẫn sửa lỗi vào ngữ cảnh hội thoại và thực hiện lại, tối đa $K$ lần (cấu hình riêng cho từng pha).

Trước khi ghi tệp, mã phải vượt qua cổng kiểm soát chất lượng theo nguyên tắc \textit{fail-fast}: (i)~kiểm tra cú pháp, (ii)~xác minh lớp giao diện bắt buộc (tuân thủ RelBench~\cite{robinson2024relbench}), và (iii)~phát hiện thư viện thiếu khai báo import. Mã không hợp lệ không được ghi vào tệp, ngăn lỗi lan truyền sang giai đoạn sau.

\subsection{Tầng 2: Gỡ lỗi nội bộ của Tác tử Vận hành}

Tầng này khác biệt về bản chất so với Tầng 1. Ở Tầng 1, tác tử thử lại toàn bộ quá trình sinh mã khi sản phẩm chưa tồn tại. Ở Tầng 2, mã đã hoàn chỉnh nhưng thất bại tại thời điểm thực thi (runtime error). Chiến lược phục hồi chuyển từ \textit{sinh lại} sang \textit{vá lỗi có hướng dẫn} (guided patching), được mô tả trong Mục~\ref{sec:operation-agent}.

Hai quyết định thiết kế trong vòng lặp gỡ lỗi của $\mathcal{A}_{ops}$:

\begin{itemize}
    \item \textbf{Khôi phục tệp gốc trước mỗi lần thử}: Mỗi vòng bắt đầu bằng khôi phục toàn bộ tệp về bản gốc. Không có bước này, bản vá sai ở vòng $i$ tạo lỗi mới ở vòng $i{+}1$ --- hiện tượng tích lũy lỗi khiến xác suất phục hồi giảm nhanh theo số vòng. Việc khôi phục biến các lần thử trở nên độc lập.
    \item \textbf{Quy chiếu lỗi đa tệp}: Lỗi runtime trong kịch bản huấn luyện có thể bắt nguồn từ tệp định nghĩa đồ thị hoặc tệp tác vụ. Hệ thống phân tích traceback để xác định tệp nguồn, và LLM sinh bản vá cho đúng tệp cần sửa.
\end{itemize}

Khi cả $N$ vòng thất bại, $\mathcal{A}_{ops}$ khôi phục tệp gốc trước khi trả quyền cho tầng điều phối --- điều kiện cần để Tầng 3 sinh lại mã từ đầu mà không bị ảnh hưởng bởi bản vá hỏng.

\subsection{Tầng 3: Phục hồi tại tầng điều phối}

Khi Tầng 1 hoặc Tầng 2 không phục hồi được, lỗi chuyển lên tầng điều phối. Hàm định tuyến phát hiện lỗi qua trạng thái toàn cục $\mathcal{S}$ và chuyển luồng đến nút xử lý lỗi, thực hiện hai bước: (1)~so sánh số lần thử với ngưỡng tối đa cho pha bị lỗi --- nếu vượt ngưỡng, pipeline kết thúc với trạng thái thất bại; (2)~xóa lỗi tích lũy, ghi nhật ký lỗi vào lịch sử (append-only), tăng bộ đếm thử lại, và chuyển luồng quay lại pha bị lỗi.

Đặc biệt, khi $\mathcal{A}_{ops}$ thất bại, tầng điều phối không quay lại pha thực thi mà quay về pha mô hình hóa --- $\mathcal{A}_{gnn}$ sinh lại kịch bản huấn luyện từ đầu thay vì thực thi lại kịch bản cũ.

\paragraph{Phản hồi có cấu trúc xuyên suốt.}
Tại mỗi tầng, ngữ cảnh lỗi được cung cấp cho LLM theo dạng có cấu trúc~\cite{shinn2023reflexion}: ở Tầng 1 qua thông điệp hệ thống trong lịch sử hội thoại, ở Tầng 2 qua traceback và lịch sử lỗi tích lũy, và ở Tầng 3 qua nhật ký lỗi trong $\mathcal{S}$ được tự động đưa vào ngữ cảnh khi tác tử được gọi lại.
